\documentclass{article}
\usepackage[utf8]{inputenc}
\usepackage{amsmath, amssymb, amsthm}
\usepackage{kotex}
\usepackage{listings}

\title{Unobservable Epistemic Mathematics (UEM): \\ A Formal System for Epistemic Limits and Margin}
\author{UEM Agent}
\date{\today}

\newtheorem{definition}{Definition}[section]
\newtheorem{axiom}{Axiom}[section]
\newtheorem{theorem}{Theorem}[section]
\newtheorem{lemma}[theorem]{Lemma}

\begin{document}

\maketitle

\begin{abstract}
This paper presents Unobservable Epistemic Mathematics (UEM), a formal mathematical system designed to model the epistemic limits of observers and the conservation of information through 'margins'. We define the object hierarchy (Scalar to Secare), the 9-dimensional epistemic coordinate system, and the Nullification Projection operator. We also provide formal proofs for core theorems using Lean 4.
\end{abstract}

\tableofcontents

\section{Introduction}
\label{sec:intro}
UEM aims to formalize the "unobservable" and "epistemic limits" within a mathematical framework. It introduces the concept of 'Margin' ($M$) to preserve information that is excluded from explicit representation ($E$). 

\section{Axioms and Object Hierarchy}
\label{sec:axioms}
The UEM object hierarchy consists of 7 layers:
\begin{itemize}
    \item Classic: Scalar, Vector, Tensor
    \item Epistemic: Sparke (\pentagram), Actyon (\bigcirc), Escalade (\leadsto), Secare (\heartsuit)
\end{itemize}

\subsection{Sparke Monoid}
We formally prove that the set of Sparkes forms a commutative monoid under the merge operation.
\begin{lstlisting}[language=Lean]
instance : AddCommMonoid (Sparke Val) := by
  apply AddCommMonoid.mk
  ...
\end{lstlisting}

\section{Nullification Projection and Dimensions}
\label{sec:nullification}
The Nullification Projection $\Pi_{null}$ maps a vector space $V = V_{keep} \oplus V_{null}$ to $V_{keep}$, while transferring $V_{null}$ information to the Margin $M$.

\section{Theorems and Formal Proofs}
\label{sec:proofs}
\subsection{Theorem P1: Solution Preservation}
We prove that $\Pi_{null}$ preserves the solution set of linear equations for the kept variables.
\begin{theorem}[P1]
Let $A: V \to W$ be a linear map such that $V_{null} \subseteq \ker A$. Then the solution set for $A(v) = b$ is isomorphic to the solution set of the projected system.
\end{theorem}
(See `formal/UEM/Theorems/P1_NullProjection.lean` for the Lean 4 proof.)

\section{Applications and Comparisons}
\label{sec:applications}
We compare UEM with classical approaches in high school algebra, calculus, and linear algebra, demonstrating how UEM's "Exclusion First" principle reduces cognitive load and computational complexity.

\end{document}